\section{Introduction}
\label{sec:intro}

Mobile apps have become the primary surface through which people manage their
daily lives, from sending messages to tracking expenses to navigating
unfamiliar cities. A mobile GUI agent is a system that operates these apps on a
user's behalf: it reads what is on the screen, decides on a tap, swipe, or text
input, and observes the result, repeating until the user's request is fulfilled.
With strong vision-language models behind them, recent agents have begun to
handle multi-step requests reliably on a fixed set of apps, opening the door to
personal assistants that act, not just chat.

The past two years have produced a wave of capable GUI
agents~\citep{cheng2024seeclick,hong2024cogagent,yang2023appagent,guo2025agentsama}.
SeeClick~\citep{cheng2024seeclick} introduced grounding-aware pretraining for
screen elements, CogAgent~\citep{hong2024cogagent} scaled a GUI-specialized
vision-language model with high-resolution support, and
AppAgent~\citep{yang2023appagent} recorded UI documentation during an offline
exploration phase to guide deployment. Benchmarks have advanced in parallel.
AndroidWorld~\citep{rawles2024androidworld} provides programmatic tasks with
reward signals derived from the device state. AndroidLab~\citep{xu2024androidlab}
adds a reproducible training framework, AndroidControl~\citep{li2024androidcontrol}
contributes a 15K-demonstration dataset over 833 apps, and
MobileWorld~\citep{mobileworld2025} extends task horizons and introduces
multi-application and MCP-augmented categories. The best systems now reach
60--80\% success rates on these benchmarks~\citep{openmobile2026,mobilerl2025}.

These headline numbers, however, are typically measured where training and test
tasks share the same template family, so they say little about how an agent
behaves on an application it has never seen. When generalization is measured
directly, the picture changes sharply. \citet{gu2026generalization} formalize
the problem as a contextual Markov decision process and build
AndroidWorld-Generalization, a benchmark that separates zero-shot generalization
to unseen task \emph{instances}, unseen \emph{templates}, and unseen
\emph{applications}. Their result is telling: online reinforcement learning
lifts a 7B vision-language agent by $+26.1\%$ on unseen instances but by only
$+8.3\%$ on unseen applications---app-level transfer is by far the hardest
regime, and it is exactly the one that matters when a user installs something
new. Crucially, the same work observes that \emph{few-shot adaptation at test
time improves performance on unseen apps}, and explicitly flags this as a
direction for future research. This paper takes up that direction.

The deployment setting constrains how it can be pursued. Synthesizing thousands
of task variants and fine-tuning per app is not viable when a user installs
several new apps in a week, and large-budget online RL is not viable when the
agent is expected to start working within minutes. The natural budget at
deployment time is on the order of a few dozen trajectories on a small
adaptation set, and an effective agent has to make that budget count. Two ideas
point at how. The first is to give the agent structured prior knowledge that
transfers across apps, so adaptation starts from something better than scratch.
The second is to update the agent's weights during adaptation, using whatever
reward signal it can extract from its own rollouts on that small set. Recent work
in single-turn reasoning shows that doing both jointly outperforms either
alone~\citep{zhang2026espl}, but whether this carries over to the multi-step,
visual, expensive-rollout GUI setting under strict unseen-app evaluation is
unanswered.

We introduce \textbf{EvoFSM-RL}, a test-time adaptation framework designed for
this budget. The agent's prior knowledge is a two-layer finite state machine: a
lower layer captures app-specific UI structure, and an upper layer captures
category-generic workflow patterns written without any app-specific token, with
a linter enforcing the separation so the upper layer is the structurally
guaranteed unit of cross-app transfer. Same-category source apps aggregate their
upper layers into a per-category abstract library $L_C$. At deployment,
adaptation runs a joint loop: a frozen reference language model proposes
structured edits to the upper-layer FSM, \emph{and} a LoRA adapter on the policy
is updated by group-relative policy optimization. The same loop runs at
source-pool pretraining and at target-app deployment, which makes ``test-time
adaptation'' a well-defined operation rather than an ad-hoc fine-tune.

We evaluate the framework at two levels of generalization. \emph{Within a
benchmark} (AndroidWorld+), where source and target apps are disjoint app sets
under a protocol that is simultaneously pool-, category-, and template-disjoint,
static category knowledge lifts near-transfer success by $+9.3$~pp over a
zero-shot baseline and online upper-layer evolution adds another $+3.7$~pp, with
a far-transfer panel acting as a built-in null control. \emph{Across benchmarks}
(AndroidWorld+ $\to$ MobileWorld), where the agent is pretrained on one
benchmark and deployed on an independently authored one, we find that static
priors which transfer within a benchmark collapse across the benchmark
boundary---a diagnosis that motivates adapting the prior on the target rather
than shipping it verbatim. The full joint adaptation, which couples upper-layer
evolution to the weight update, is the central experiment this framework sets
up.

\medskip
\noindent\textbf{Contributions.} This paper makes the following contributions.
\begin{itemize}
    \item \textbf{Test-time adaptation for unseen apps.} We take on the hardest
    and most deployment-relevant generalization regime identified
    by~\citet{gu2026generalization}---transfer to unseen applications---and
    develop the few-shot test-time adaptation it flags as promising into a full
    framework, EvoFSM-RL, that adapts a GUI agent to a new app on a small
    adaptation budget rather than relying on per-app synthesis or large-budget
    retraining.

    \item \textbf{A joint prompt-and-weight adaptation algorithm.} At the core is
    a deployment-time loop (our B4) that co-evolves two channels at once: it
    mutates a transferable upper-layer FSM via a frozen reference language model
    \emph{and} updates a LoRA adapter via group-relative policy optimization with
    within-(FSM, task) baselines. The same loop runs at source-pool pretraining
    and at target-app adaptation, with the category library $L_C$ mediating
    between them, so prompt evolution and weight adaptation are co-adapted rather
    than treated as separate stages.

    \item \textbf{A multi-level benchmark partition for measuring
    generalization.} We partition mobile-agent benchmark data to expose model
    ability at distinct generalization levels: within a benchmark, a protocol
    that is simultaneously pool-, category-, and template-disjoint; and across
    benchmarks, a category-level near/far transfer tiering between an
    AndroidWorld+ training source and an independently authored MobileWorld
    deployment target. This extends the instance/template/app regimes
    of~\citet{gu2026generalization} with category-level tiers and a
    cross-benchmark axis.
\end{itemize}
